\section{Alternating Series and Absolute Convergence}\label{sec:alt_series}

The series convergence tests we have used require that the underlying sequence $\{a_n\}$ be a positive sequence. (We can relax this with \autoref{thm:series_behavior} and state that there must be an $N>0$ such that $a_n>0$ for all $n>N$; that is, $\{a_n\}$ is positive for all but a finite number of values of $n$.)

In this section we explore series whose summation includes negative terms. We start with a very specific form of series, where the terms of the summation alternate between being positive and negative.

\begin{definition}[Alternating Series]\label{def:alt_series}%
Let $\{b_n\}$ be a positive sequence. An \textbf{alternating series} is a series of either the form
\index{series!alternating}
\[
\sum_{n=1}^\infty (-1)^nb_n\qquad \text{or}\qquad \sum_{n=1}^\infty (-1)^{n+1}b_n.
\]
\end{definition}

We want to think that an alternating sequence $\{a_n\}$ is related to a positive sequence $\{b_n\}$ by $a_n=(-1)^n b_n$.

Recall that the terms of Harmonic Series come from the Harmonic Sequence $\{b_n\} = \{\frac1n\}$. An important alternating series is the \textbf{Alternating Harmonic Series}:
\[
\sum_{n=1}^\infty (-1)^{n+1}\frac1n
= 1-\frac12+\frac13-\frac14+\frac15-\frac16+\dotsb
\]

Geometric Series can also be alternating series when $r<0$. For instance, if $r=-1/2$, the geometric series is
\[
\sum_{n=0}^\infty \left(\frac{-1}{2}\right)^n
= 1-\frac12+\frac14-\frac18+\frac1{16}-\frac1{32}+\dotsb
\]

\autoref{thm:geom_series} states that geometric series converge when $\abs r<1$ and gives the sum: $\ds \sum_{n=0}^\infty r^n = \frac1{1-r}$. When $r=-1/2$ as above, we find
\[
\sum_{n=0}^\infty\left(\frac{-1}{2}\right)^n=\frac1{1-(-1/2)}=\frac 1{3/2}=\frac23.
\]

A powerful convergence theorem exists for other alternating series that meet a few conditions.

\begin{theorem}[Alternating Series Test]\label{thm:alt_series_test}%
Let $\{b_n\}$ be a positive, decreasing sequence where $\ds\lim_{n\to\infty}b_n=0$. Then
\index{series!Alternating Series Test}\index{Alternating Series Test!for series}\index{convergence!Alternating Series Test}\index{divergence!Alternating Series Test}
\[
\sum_{n=1}^\infty (-1)^{n}b_n \qquad \text{and}\qquad \sum_{n=1}^\infty (-1)^{n+1}b_n
\]
converge.
\end{theorem}

The basic idea behind \autoref{thm:alt_series_test} is illustrated in \autoref{fig:alt_series_converge}. A positive, decreasing sequence $\{b_n\}$ is shown along with the partial sums
\[S_n = \sum_{i=1}^n(-1)^{i+1}b_i =b_1-b_2+b_3-b_4+\dotsb+(-1)^{n+1}b_n.\]
Because $\{b_n\}$ is decreasing, the amount by which $S_n$ bounces up and down decreases. Moreover, the odd terms of $S_n$ form a decreasing, bounded sequence, while the even terms of $S_n$ form an increasing, bounded sequence. Since bounded, monotonic sequences converge (see \autoref{thm:monotonic_converge}) and the terms of $\{b_n\}$ approach 0, we will show below that the odd and even terms of $S_n$ converge to the same common limit $L$, the sum of the series.

\mtable{Illustrating convergence with the Alternating Series Test.}{fig:alt_series_converge}{\begin{tikzpicture}[alt={The sequence bₙ is decreasing from 1 toward 0.  The partial sums Sₙ alternate on which side of the dashed line y=L they are.}]
\begin{axis}[width=1.16\marginparwidth,tick label style={font=\scriptsize},
axis y line=middle,axis x line=middle,name=myplot,axis on top,xtick={2,4,6,8,10},
ymin=-.25,ymax=1.25,xmin=-1,xmax=11,legend columns=-1,legend style={at={(.5,0)},anchor=north}]
\draw [thick, dashed] (axis cs: -.1,.614) node [left] {\scriptsize $L$} -- (axis cs: 10.25,.614);
\addplot [only marks,draw={\colorone},fill={\colorone},mark size={1.5pt},domain=1:10,samples=10] {2/(x+1)};
\addlegendentry{\scriptsize$b_n$\quad\mbox{}}
\addplot [only marks,draw={\colortwo},fill={\colortwo},mark size={1.5pt},domain=1:10,samples=10]coordinates {(1., 1.)(2., 0.333333)(3., 0.833333)(4., 0.433333)(5.,0.766667)(6., 0.480952)(7., 0.730952)(8., 0.50873)(9.,0.70873)(10., 0.526912)};
\addlegendentry{\scriptsize$S_n$}
\end{axis}
\node [right] at (myplot.right of origin) {\scriptsize $n$};
\node [above] at (myplot.above origin) {\scriptsize $y$};
\end{tikzpicture}}

\begin{proof}
Because $\{b_n\}$ is a decreasing sequence, we have $b_n-b_{n+1}\geq 0$. 
We will consider the even and odd partial sums separately. First consider the even partial sums.
\begin{align*}
S_2&=b_1-b_2\geq 0  &\text{since } b_2&\leq b_1\\
S_4&=b_1-b_2+b_3-b_4=S_2+b_3-b_4\geq S_2 &\text{since } b_3- b_4&\geq 0\\
S_6&=S_4+b_5-b_6\geq S_4 &\text{since } b_5- b_6&\geq 0\\
\vdots \\
S_{2n}&=S_{2n-2}+b_{2n-1}-b_{2n}\geq S_{2n-2}  &\text{since } b_{2n-1}- b_{2n}&\geq 0
\end{align*}
We now have
\[0\leq S_2\leq S_4\leq S_6\leq  \dotsb \leq S_{2n}\leq \dotsb\]
so $\{S_{2n}\}$ is an increasing sequence. 
But we can also write 
\begin{align*}
S_{2n}&=b_1-b_2+b_3-b_4+b_5-\dotsb -b_{2n-2}+b_{2n-1}-b_{2n}\\
&=b_1-(b_2-b_3)-(b_4-b_5)-\dotsb -(b_{2n-2}-b_{2n-1})-b_{2n}
\end{align*}
Each term in parentheses is positive and $b_{2n}$ is positive so we have $S_{2n}\leq b_1$ for all $n$. We now have the sequence of even partial sums, $\{S_{2n}\}$, is increasing and bounded above so by \autoref{thm:monotonic_converge}
 $\{S_{2n}\}$ converges. Since we know it converges, we will assume it's limit is $L$ or
\[\lim_{n\to \infty}S_{2n}=L\]
Next we determine the limit of the sequence of odd partial sums.
\begin{align*}
\lim_{n\to \infty}S_{2n+1}&=\lim_{n\to \infty}(S_{2n}+b_{2n+1})\\
&=\lim_{n\to \infty}S_{2n}+\lim_{n\to \infty}b_{2n+1}\\
&=L+0\\
&=L
\end{align*}
Both the even and odd partial sums converge to $L$ so we have $\ds \lim_{n\to \infty}S_n=L$, which means the series is convergent.
\end{proof}

\youtubeVideo{aOiZvfFAMW8}{Alternating Series --- Another Example 4}

\begin{example}[Applying the Alternating Series Test]\label{ex_alt1}%
Determine if the Alternating Series Test applies to each of the following series.
\[
 \text{1. }\sum_{n=1}^\infty (-1)^{n+1}\frac1n\qquad
 \text{2. }\sum_{n=2}^\infty (-1)^n\frac{\ln n}{n}\qquad
 \text{3. }\sum_{n=1}^\infty (-1)^{n+1}\frac{\abs{\sin n}}{n^2}
\]
\solution
\begin{enumerate}
	\item This is the Alternating Harmonic Series as seen previously. The underlying sequence is $\{b_n\} = \{1/n\}$, which is positive, decreasing, and approaches 0 as $n\to\infty$. Therefore we can apply the Alternating Series Test and conclude this series converges. 
	
	While the test does not state what the series converges to, we will see later that $\ds \sum_{n=1}^\infty (-1)^{n+1}\frac1n=\ln2.$
	
	\item		The underlying sequence is $\{b_n\} = \{(\ln n)/n\}$. This is positive for $n\geq 2$ and $\ds \lim_{n\to \infty} \frac{\ln n}{n}=\lim_{n\to \infty}\frac{1}{n}=0$ (use L'Hôpital's Rule). However, the sequence is not decreasing for all $n$. It is straightforward to compute $b_1\approx0.347$, $b_2\approx 0.366$, and $b_3\approx 0.347$: the sequence is increasing for at least the first 2 terms. 
	
	We do not immediately conclude that we cannot apply the Alternating Series Test. Rather, consider the long-term behavior of $\{b_n\}$. Treating $b_n=b(n)$ as a continuous function of $n$ defined on $[2,\infty)$, we can take its derivative:
	\[b\primeskip'(n) = \frac{1-\ln n}{n^2}.\]
	The derivative is negative for all $n\geq 3$ (actually, for all $n>e$), meaning $b(n)=b_n$ is decreasing on $[3,\infty)$. We can apply the Alternating Series Test to the series when we start with $n=3$ and conclude that $\ds \sum_{n=3}^\infty(-1)^n\frac{\ln n}{n}$ converges; adding the terms with $n=2$ does not change the convergence (i.e., we apply \autoref{thm:series_behavior}).
	
	The important lesson here is that as before, if a series fails to meet the criteria of the Alternating Series Test on only a finite number of terms, we can still apply the test.
	
	\item  The underlying sequence is $\{b_n\} = \{\abs{\sin n}/n^2\}$. This sequence is positive and approaches $0$ as $n\to\infty$. However, it is not a decreasing sequence; the value of $\abs{\sin n}$ oscillates between $0$ and $1$ as $n\to\infty$. We cannot remove a finite number of terms to make $\{b_n\}$ decreasing, therefore we cannot apply the Alternating Series Test.
	
	Keep in mind that this does not mean we conclude the series diverges; in fact, it does converge. We are just unable to conclude this based on \autoref{thm:alt_series_test}.
\end{enumerate}
\end{example}

One of the famous results of mathematics is that the Harmonic Series,\vspace{-.5\baselineskip} $\ds \sum_{n=1}^\infty \frac1n$ diverges, yet the Alternating Harmonic Series, $\ds \sum_{n=1}^\infty (-1)^{n+1}\frac1n$, converges. The notion that alternating the signs of the terms in a series can make a series converge leads us to the following definitions.\index{Alternating Harmonic Series}

\begin{definition}[Absolute and Conditional Convergence]\label{def:abs_converge}%
\index{convergence!absolute}\index{convergence!conditional}\index{series!absolute convergence}\index{series!conditional convergence}%
\mbox{}\\[-2\baselineskip]\parbox[t]{\linewidth}{%
\begin{enumerate}
	\item A series $\ds \sum_{n=1}^\infty a_n$ \textbf{converges absolutely} if $\ds \sum_{n=1}^\infty \abs{a_n}$ converges.
	\item A series $\ds \sum_{n=1}^\infty a_n$ \textbf{converges conditionally} if $\ds \sum_{n=1}^\infty a_n$ converges but $\ds \sum_{n=1}^\infty \abs{a_n}$ diverges.
\end{enumerate}}
\end{definition}

\mnote{\textbf{Note:} In \autoref{def:abs_converge}, $\ds \sum_{n=1}^\infty a_n$ is not necessarily an alternating series; it just may have some negative terms.}

Thus we say the Alternating Harmonic Series converges conditionally.

\begin{example}[Determining absolute and conditional convergence.]\label{ex_alt_series2}%
Determine if the following series converge absolutely, conditionally, or diverge.\\
\begin{minipage}{.5\textwidth}
 \begin{enumerate}
  \item $\ds\sum_{n=1}^\infty (-1)^n\frac{n+3}{n^2+2n+5}$
%  \item $\ds\sum_{n=1}^\infty (-1)^n\frac{n^2+2n+5}{2^n}$ % ratio test is next section
 \end{enumerate}
\end{minipage}%
\begin{minipage}{.5\textwidth}
 \begin{enumerate}\setcounter{enumi}{1}
  \item $\ds\sum_{n=3}^\infty (-1)^n\frac{3n-3}{5n-10}$
 \end{enumerate}
\end{minipage}
\solution
\begin{enumerate}
	\item We can show the series
	\[\sum_{n=1}^\infty \abs{(-1)^n\frac{n+3}{n^2+2n+5}}= \sum_{n=1}^\infty \frac{n+3}{n^2+2n+5}\]
	diverges using the Limit Comparison Test, comparing with $1/n$. 
	
	The sequence $\{\dfrac{n+3}{n^2+2n+5}\}$ is monotonically decreasing, so that the series $\ds \sum_{n=1}^\infty (-1)^n\frac{n+3}{n^2+2n+5}$ converges using the Alternating Series Test; we conclude it converges conditionally.
	
%	\item	We can show the series
%	\[\sum_{n=1}^\infty \abs{(-1)^n\frac{n^2+2n+5}{2^n}}=\sum_{n=1}^\infty \frac{n^2+2n+5}{2^n}\]
%	converges using the Ratio Test. 
%	
%	Therefore we conclude $\ds \sum_{n=1}^\infty (-1)^n\frac{n^2+2n+5}{2^n}$ converges absolutely.
	
	\item	The series
	\[\sum_{n=3}^\infty \abs{(-1)^n\frac{3n-3}{5n-10}} = \sum_{n=3}^\infty \frac{3n-3}{5n-10}\]
	diverges using the Test for Divergence, so it does not converge absolutely. 
	
	The series $\ds \sum_{n=3}^\infty (-1)^n\frac{3n-3}{5n-10}$ fails the conditions of the Alternating Series Test as $(3n-3)/(5n-10)$ does not approach $0$ as $n\to\infty$. We can state further that this series diverges; as $n\to\infty$, the series effectively adds and subtracts $3/5$ over and over. This causes the sequence of partial sums to oscillate and not converge.
	
	Therefore the series $\ds \sum_{n=1}^\infty (-1)^n\frac{3n-3}{5n-10}$ diverges.
\end{enumerate}
\end{example}

Knowing that a series converges absolutely allows us to make two important statements, given in the following theorem. The first is that absolute convergence is  ``stronger'' than regular convergence.\vspace{-.5\baselineskip}
That is, just because {\small$\ds \sum_{n=1}^\infty a_n$} converges, we cannot conclude that {\small$\ds \sum_{n=1}^\infty \abs{a_n}$} will converge, but knowing a series converges absolutely tells us that {\small$\ds \sum_{n=1}^\infty a_n$} will converge. 

\begin{theorem}[Absolute Convergence Theorem]\label{thm:abs_convergence}%
Let $\ds \sum_{n=1}^\infty a_n$ be a series that converges absolutely.
\index{convergence!absolute}\index{Absolute Convergence Theorem}\index{series!Absolute Convergence Theorem}\index{rearrangements of series}\index{series!rearrangements}
\begin{enumerate}
	\item $\ds \sum_{n=1}^\infty a_n$ converges.
	\item	Let $\{b_n\}$ be any rearrangement of the sequence $\{a_n\}$. Then 
	\[\sum_{n=1}^\infty b_n = \sum_{n=1}^\infty a_n.\]
\end{enumerate}
\end{theorem}

One reason this is important is that our convergence tests all require that the underlying sequence of terms be positive. By taking the absolute value of the terms of a series where not all terms are positive, we are often able to apply an appropriate test and determine absolute convergence. This, in turn, determines that the series we are given also converges.

The second statement relates to \textbf{rearrangements}\index{rearrangements of series}\index{series!rearrangements} of series. When dealing with a finite set of numbers, the sum of the numbers does not depend on the order which they are added. (So $1+2+3 = 3+1+2$.) One may be surprised to find out that when dealing with an infinite set of numbers, the same statement does not always hold true: some infinite lists of numbers may be rearranged in different orders to achieve different sums. The theorem states that the terms of an absolutely convergent series can be rearranged in any way without affecting the sum.

% todo Tim do we want to have an absolutely convergent example to refer back to?
% the former part 2 was removed because it used the ratio test
%In \autoref{ex_alt_series2}, we determined the series in part 2 converges absolutely. \autoref{thm:abs_convergence} tells us the series converges (which we could also determine using the Alternating Series Test).

The theorem states that rearranging the terms of an absolutely convergent series does not affect its sum. This implies that perhaps the sum of a conditionally convergent series can change based on the arrangement of terms. Indeed, it can. The Riemann Rearrangement Theorem (named after Bernhard Riemann) states that any conditionally convergent series can have its terms rearranged so that the sum is any desired value or infinity.

Before we consider an example, we state the following theorem that illustrates how the alternating structure of an alternating series is a powerful tool when approximating the sum of a convergent series.

\begin{theorem}[The Alternating Series Approximation Theorem]\label{thm:alt_series_approx}%
Let $\{b_n\}$ be a sequence that satisfies the hypotheses of the Alternating Series Test, and let $S_n$ and $L$ be the $n^{\text{th}}$ partial sum and sum, respectively, of either $\ds \sum_{n=1}^\infty (-1)^{n}b_n$ or $\ds \sum_{n=1}^\infty (-1)^{n+1}b_n$. Then
\index{series!alternating!Approximation Theorem}
\begin{enumerate}
	\item $\abs{S_n-L} < b_{n+1}$, and
	\item	$L$ is between $S_n$ and $S_{n+1}$.
\end{enumerate}
\end{theorem}

Part 1 of \autoref{thm:alt_series_approx} states that the $n^{\text{th}}$ partial sum of a convergent alternating series will be within $b_{n+1}$ of its total sum. Consider the alternating series we looked at before the statement of the theorem, $\ds \sum_{n=1}^\infty \frac{(-1)^{n+1}}{n^2}$. Since $b_{14} = 1/14^2 \approx 0.0051$, we know that $S_{13}$ is within $0.0051$ of the total sum. 

Moreover, Part 2 of the theorem states that since $S_{13} \approx 0.8252$ and $S_{14}\approx 0.8201$, we know the sum $L$ lies between $0.8201$ and $0.8252$. One use of this is the knowledge that $S_{14}$ is accurate to two places after the decimal.

Some alternating series converge slowly. In \autoref{ex_alt1} we determined the series $\ds\sum_{n=2}^\infty (-1)^{n+1}\frac{\ln n}{n}$ converged. With $n=1001$, we find $(\ln n)/n \approx 0.0069$, meaning that $S_{1000} \approx 0.1633$ is accurate to one, maybe two, places after the decimal. Since $S_{1001} \approx 0.1564$, we know the sum $L$ is $0.1564\leq L\leq0.1633$.

\begin{example}[Approximating the sums of convergent alternating series]\label{ex_alt_series_approx}%
Approximate the sum of the following series, accurate to within $0.001$.
\[
 \text{1. }\sum_{n=1}^\infty (-1)^{n+1}\frac{1}{n^3}\qquad
 \text{2. }\sum_{n=1}^\infty (-1)^{n+1}\frac{\ln n}{n}.
\]
\solution
\begin{enumerate}
	\item  Using \autoref{thm:alt_series_approx}, we want to find $n$ where $1/n^3 \le 0.001$:
	\begin{align*}
	\frac1{n^3} &\leq 0.001=\frac{1}{1000} \\
	n^3 &\geq 1000\\
	n &\geq \sqrt[3]{1000}\\
	n &\geq 10.
	\end{align*}
	Let $L$ be the sum of this series. By Part 1 of the theorem, $\abs{S_9-L}<b_{10} = 1/1000$. We can compute $S_9=0.902116$, which our theorem states is within $0.001$ of the total sum. 
	
	We can use Part 2 of the theorem to obtain an even more accurate result. As we know the 10\textsuperscript{th} term of the series is $-1/1000$, we can easily compute $S_{10} = 0.901116$. Part 2 of the theorem states that $L$ is between $S_9$ and $S_{10}$, so $0.901116 <L<0.902116$.
	
	\item		We want to find $n$ where $(\ln n)/n \le 0.001$. We start by solving $(\ln n)/n = 0.001$ for $n$. This cannot be solved algebraically, so we will use Newton's Method to approximate a solution. 
	
	Let $f(x) = \ln(x)/x-0.001$; we want to know where $f(x) = 0$. We make a guess that $x$ must be ``large,'' so our initial guess will be $x_1=1000$. Recall how Newton's Method works: given an approximate solution $x_n$, our next approximation $x_{n+1}$ is given by
	\[x_{n+1} = x_n - \frac{f(x_n)}{\fp(x_n)}.\]
	We find $\fp(x) = \bigl(1-\ln(x)\bigr)/x^2$. This gives
	\begin{align*}
	x_2 &= 1000 - \frac{\ln(1000)/1000-0.001}{\bigl(1-\ln(1000)\bigr)/1000^2} \\
			&= 2000.
	\end{align*}
	Using a computer, we find that Newton's Method seems to converge to a solution $x=9118.01$ after 8 iterations. Taking the next integer higher, we have $n=9119$, where $\ln(9119)/9119 =0.000999903<0.001$.
	
	Again using a computer, we find $S_{9118} = -0.160369$. Part 1 of the theorem states that this is within $0.001$ of the actual sum $L$. Already knowing the 9,119\textsuperscript{th} term, we can compute $S_{9119} = -0.159369$, meaning
	\[-0.160369 < L < -0.159369.\]
	
\end{enumerate}
Notice how the first series converged quite quickly, where we needed only 10 terms to reach the desired accuracy, whereas the second series took over 9,000 terms.
\end{example}

We now consider the Alternating Harmonic Series once more. We have stated that 
\[
\sum_{n=1}^\infty (-1)^{n+1}\frac1n
= 1-\frac12+\frac13-\frac14+\frac15-\frac16+\frac17\dotsb = \ln 2,
\]
(see \autoref{ex_alt1}). 

Consider the rearrangement where every positive term is followed by two negative terms:
\[
1-\frac12-\frac14+\frac13-\frac16-\frac18+\frac15-\frac1{10}-\frac1{12}\dotsb
\]
(Convince yourself that these are exactly the same numbers as appear in the Alternating Harmonic Series, just in a different order.) Now group some terms and simplify:
\begin{align*}
\left(1-\frac12\right)-\frac14+\left(\frac13-\frac16\right)-\frac18
+\left(\frac15-\frac1{10}\right)-\frac1{12}+\dotsb &= \\
\frac12-\frac14+\frac16-\frac18+\frac1{10}-\frac{1}{12}+\dotsb &= \\
\frac12\left(1-\frac12+\frac13-\frac14+\frac15-\frac16+\dotsb\right) & = \frac12\ln 2.
\end{align*}

By rearranging the terms of the series, we have arrived at a different sum. (One could \emph{try} to argue that the Alternating Harmonic Series does not actually converge to $\ln 2$, because rearranging the terms of the series \emph{shouldn't} change the sum. However, the Alternating Series Test proves this series converges to $L$, for some number $L$, and if the rearrangement does not change the sum, then $L = L/2$, implying $L=0$. But the Alternating Series Approximation Theorem quickly shows that $L>0$. The only conclusion is that the rearrangement \emph{did} change the sum.) This is an incredible result.\bigskip


We mentioned earlier that the Integral Test did not work well with series containing factorial terms. The next section introduces the Ratio Test, which does handle such series well. We also introduce the Root Test, which is good for series where each term is raised to a power.

\printexercises{exercises/08-05-exercises}
